\documentclass[11pt]{amsart}
\usepackage[T1]{fontenc}
\usepackage{lmodern}
\usepackage{amsmath,amssymb,amsthm,mathtools}
\usepackage{microtype}
\usepackage[margin=1.05in]{geometry}
\usepackage{xurl}
\usepackage[hidelinks]{hyperref}
\newcommand{\Z}{\mathbb Z}
\newcommand{\F}{\mathbb F}
\DeclareMathOperator{\GL}{GL}
\DeclareMathOperator{\St}{St}
\DeclareMathOperator{\id}{id}
\DeclareMathOperator{\im}{im}
\DeclareMathOperator{\coker}{coker}
\DeclareMathOperator{\Hom}{Hom}
\DeclareMathOperator{\Ext}{Ext}
\DeclareMathOperator{\Tor}{Tor}
\newtheorem{theorem}{Theorem}[section]
\newtheorem{lemma}[theorem]{Lemma}
\newtheorem{proposition}[theorem]{Proposition}
\newtheorem{corollary}[theorem]{Corollary}
\theoremstyle{definition}
\newtheorem{definition}[theorem]{Definition}
\theoremstyle{remark}
\newtheorem{remark}[theorem]{Remark}
\numberwithin{equation}{section}

% Display numbers for outer $$ ... $$ delimiters.
\newcommand{\eqnum}[1][]{%
  \refstepcounter{equation}%
  \if\relax\detokenize{#1}\relax\else\label{#1}\fi
  \leqno\hbox{\normalfont(\theequation)}}
\newcounter{displaytag}

% Row tags for a single $$ ... aligned ... $$ display.
\newcommand{\eqrownum}[1][]{%
  \hbox{\refstepcounter{equation}%
    \if\relax\detokenize{#1}\relax\else\label{#1}\fi
    \normalfont(\theequation)}}
\makeatletter
\newcommand{\eqrowtag}[2]{%
  \hbox{\refstepcounter{displaytag}%
    \def\@currentlabel{#1}%
    \label{#2}%
    \normalfont(#1)}}
\makeatother

\emergencystretch=1.5em
\title{General linear and Steinberg groups over the Leavitt algebra $L_{\mathbb F_2}(1,2)$}
\author{Huynh Viet Khanh}
\hypersetup{pdfauthor={Huynh Viet Khanh},pdftitle={General linear and Steinberg groups over the Leavitt algebra L\_F2(1,2)}}
\makeatletter
\renewcommand{\@setauthors}{%
  \begin{center}
  \normalfont
  {\large\authors\par}
  \medskip
  {\small Department of Mathematics and Informatics, HCMC University of Education\\
  280 An Duong Vuong Str., Ho Chi Minh City, Vietnam\\
  \href{mailto:khanhhv@hcmue.edu.vn}{\texttt{khanhhv@hcmue.edu.vn}}\par}
  \end{center}%
}
\makeatother
\date{}
\subjclass[2020]{Primary 20J05, 16S88; Secondary 19C20, 16U60, 20F05}
\keywords{Leavitt algebras, groups of units, group homology, acyclic groups, Steinberg groups, finite presentations}
\begin{document}
\begin{abstract}
Let $R=L_{\F_2}(1,2)$. We prove that $\GL_r(R)$ is integrally acyclic for every $r\geq1$ and that the canonical map $\St_r(R)\to\GL_r(R)$ is an isomorphism for every $r\geq3$. The homology calculation combines simultaneous extensions of ordered frames with scalar actions of the multiplicative groups of finite fields on their stabilizers. The presentation associated with the same frame complex defines a surjective section of the Steinberg map. An explicit finite presentation of $R^\times$ then follows from the theorem of Krsti\'c and McCool. We formulate separate criteria for acyclicity and for the Steinberg comparison over other rings.
\end{abstract}
\maketitle

\section{Introduction}

Let $R=L_{\F_2}(1,2)$ be the binary Leavitt algebra and let $G=R^\times$. We use the canonical $\F_2$-linear involution on $R$, interchanging the defining generators $e,f$ with $e^*,f^*$. The defining relations determine isomorphisms $R\cong M_r(R)$ for every $r\geq1$, expressed by coordinates indexed by finite binary leaf sets. These are instances of the isomorphisms with matrix rings studied by Abrams, \'{A}nh and Pardo \cite{AAP08}; we use the explicit coordinate maps below. We prove that $\GL_r(R)$ is integrally acyclic for every $r\geq1$ and that the ordinary Steinberg map is an isomorphism for every $r\geq3$. Together with the theorem of Krsti\'c and McCool on finite presentability, the latter result leads to an explicit finite presentation of $G$.

The relation between Leavitt algebras and Higman--Thompson groups has been studied by Pardo \cite{Pardo11} and, for unfolding trees, by Gorazd \cite{Gorazd26}. Freeland introduced the subgroup generated by scalar matrix groups associated with finite leaf sets and asked whether this subgroup contains every unit \cite[Definitions~4.3.9--4.3.10 and pp.~173--174]{Freeland19}. Khanh and Thanh consider finite presentability through unstable $K_2$ \cite[\S8]{KT26}. More generally, Preusser describes generators for the general linear groupoid of a Leavitt path algebra associated with a finite graph \cite[Theorem~4.13]{Preusser26}.

Szymik and Wahl prove integral acyclicity of Thompson's group $V$ as part of their calculation of Higman--Thompson homology \cite{SW19}. Li proves acyclicity results for topological full groups of ample groupoids \cite[Corollary~D]{Li25}, and Palmer and Wu treat labelled Thompson groups and twisted Brin--Thompson groups \cite[Theorems~A and~C]{PW25}. The standard copy of $V$ in $G$ is proper. Indeed, $u=1+ef^*$ satisfies $u^{-1}=u$ and $u^*=1+fe^*\ne u$, whereas leaf permutations are unitary. The cited results concern Thompson groups and topological full groups; the present calculation concerns the full group of units of $R$.

The stable input is the calculation of Ara, Brustenga and Corti\~nas \cite[Theorem~7.6]{ABC09}: $K(R)$ is the homotopy cofiber of multiplication by $-1$ on $K(\F_2)$. Thus the stable general linear group has zero positive integral homology. To pass to finite rank, we first show that standard stabilization induces zero on positive homology. An explicit compression has square conjugate to itself; the stable calculation then makes its homology action zero. The required further input is the existence of common extensions of ordered frames over $R$.

Complexes of completable frames and their stabilizer spectral sequences are standard tools in homological stability; see van der Kallen \cite{vdK80}, Suslin \cite{Suslin84}, Nesterenko--Suslin \cite[\S2]{NS90}, and Randal-Williams--Wahl \cite{RWW17}. Our complex uses right modules and fixes the formal size of a complement, since free modules over $R$ need not have a unique rank. Its required homology vanishing is proved directly. A simultaneous refinement of the usual Leavitt reduction argument \cite[Theorem~2.2.11]{AAS17} supplies one right multiplier for a finite family and explicit nonzero complements. It produces a single vector extending every frame in a prescribed finite family.

The scalar calculation also has established antecedents. The weights of the multiplicative group of a finite field and their degree bound are due to Quillen \cite[\S11, Lemmas~15--16]{Quillen72}. Knudson combines these weights with the unipotent subgroups of simplex stabilizers and the spectral sequence of the group action \cite[Proposition~2.1 and Theorem~3.1]{Knudson02}. The arguments of Nesterenko--Suslin \cite[Proposition~1.10 and Theorem~1.11]{NS90} and Schlichting \cite[\S2]{Schlichting17} for affine groups provide related scalar methods. Here the finite fields embedded in $R$ need not be central. Right rescaling of all frame vectors nevertheless defines a strict $\GL_r(R)$-equivariant action, and inverse left scaling on the additive normal subgroup $U$ of each stabilizer commutes with the conjugation action of its subgroup of block diagonal matrices. Keeping this action on the entire complex permits an integral induction over all positive homological degrees.

For the ordinary Steinberg map, we use the presentation associated with an action on a simply connected complex, as in Brown \cite[Theorem~1 and \S3]{Brown84}. We verify the vertex, edge and triangle lifts in the Steinberg group itself. The resulting surjective section proves injectivity at rank four. Voronetsky's theorem on refinement of idempotents extends the conclusion to all the stated ranks \cite[\S4, Proposition~1]{Voronetsky20}. Finite presentability follows from this comparison and \cite[Theorem~3]{KM99}; the presentation in Section~\ref{sec:finite-presentation} is an explicit specialization of \cite[\S3]{KM99}.

The acyclicity and Steinberg criteria are stated separately, with their geometric and algebraic hypotheses. The main conclusions appear in Theorems~\ref{thm:acyclicity}, \ref{thm:steinberg-isomorphism}, and~\ref{thm:finite-presentation}. No assertion is made about the Steinberg kernel at rank two. Proposition~\ref{prop:explicit-presentation} records the finite presentation, its coefficient bounds, and its generator evaluation.

\section{An acyclicity criterion}
\label{sec:acyclicity-criterion}

All rings in this section are nonzero unital $\F_2$-algebras. Modules are right modules, and matrices act on columns from the left. Homology has integral coefficients unless another coefficient module is specified.

We use the usual construction by completable ordered frames; compare the special unimodular frames of \cite[\S2]{NS90} and \cite[\S2]{Suslin84}. We specify the complement size and truncate at the chosen formal rank. These conventions are needed over rings whose free modules can have different finite bases.

\begin{definition}\label{def:ordered-frames}
Let $A$ be such an algebra and write $\Gamma_r=\GL_r(A)$. The ordered frame semisimplicial set $X_r(A)$ has as its $p$-simplices the tuples $(v_1,\ldots,v_k)$, where $k=p+1\leq r$, admitting a decomposition
$$
A^r=v_1A\oplus\cdots\oplus v_kA\oplus C,\qquad\text{with }C\cong A^{r-k}. \eqnum[eq:frame-decomposition]
$$
Each map $A\to v_iA$, $a\mapsto v_i a$, must be an isomorphism. Since $A\neq0$, the condition $C\cong A^{r-k}$ implies that $C\neq0$ when $k<r$ and $C=0$ when $k=r$. The simplex records the ordered vectors and the formal integer $r-k$; the complement and its basis are not recorded. Faces delete vectors, and there are no simplices in dimensions $p\geq r$.
\end{definition}

Deleting a vector adds its summand to a witness complement, so the faces in Definition~\ref{def:ordered-frames} are well defined. The distinction between a complement and its formal size is useful when the algebra has isomorphic free modules of different finite ranks.

\begin{theorem}\label{thm:acyclicity-criterion}
Suppose that $A^2\cong A$ as right $A$-modules. Assume also that:
\begin{enumerate}
\item for every $n>0$, the standard inclusion $\Gamma_{n+2}\to\Gamma_{n+3}$, $g\mapsto\operatorname{diag}(g,1)$, induces zero on $H_n(-,\Z)$;
\item $\widetilde H_i(X_r(A),\Z)=0$ for $0\leq i\leq r-3$ and $r\geq4$.
\end{enumerate}
Then $H_n(A^\times,\Z)=0$ for every $n>0$.
\end{theorem}

The following lemma is the form of Quillen's weight calculation for finite fields \cite[\S11, Lemmas~15--16]{Quillen72} needed in characteristic two. It includes the consequences for integral homology and homology with coefficients. We include the proof to specify the scalar action and degree bound.

\begin{lemma}\label{lem:scalar-homology}
Let $m\geq1$, let $F=\F_{2^m}$ and $C=F^\times$, and regard an $F$-vector space $V$ as a discrete additive group. For $1\leq j<m$, the group $H_j(V,\Z)$ is annihilated by $2$, and the norm $N_C=\sum_{c\in C}[c]\in\Z[C]$ acts as zero on it. If a group $H$ acts $F$-linearly on $V$, then $N_C$ acts as zero on $H_i(H,H_j(V,\Z))$ for every $i\geq0$. These coefficient homology groups and all their $C$-stable subquotients have zero $C$-coinvariants. The same conclusions hold for inverse scalar multiplication.
\end{lemma}

\begin{proof}
Suppose first that $V$ is finite dimensional over $F$. Its additive group is a finite product of copies of $C_2$. The integral cellular chain complex of $BC_2=\mathbb{RP}^{\infty}$ has one copy of $\Z$ in each nonnegative degree, with differential $2$ in positive even degrees and zero in odd degrees. It is the direct sum of $\Z[0]$ and two-term complexes $\Z\xrightarrow{2}\Z$ in degrees $(2a,2a-1)$, $a\geq1$. Multiplication by $2$ on each two-term complex is null-homotopic: take the identity map from degree $2a-1$ to degree $2a$ as the homotopy.

The product chain complex is the tensor product of these complexes. Every summand other than $\Z[0]$ contains a two-term factor. Applying the preceding homotopy to one such factor, with its tensor sign, makes multiplication by $2$ null-homotopic on that summand; the cross terms cancel. Consequently $2H_j(V,\Z)=0$ for $j>0$. The coefficient sequence $\Z\xrightarrow{2}\Z\to\F_2$ therefore induces a natural injection
$$
H_j(V,\Z)\longrightarrow H_j(V,\F_2)\qquad\text{for }j>0, \eqnum[eq:integral-mod-two]
$$
which commutes with every automorphism of $V$.

The mod-two cohomology of $V$ is the polynomial algebra on $H^1(V,\F_2)=\Hom_{\F_2}(V,\F_2)$. Indeed, $H^*(C_2,\F_2)=\F_2[t]$ with $|t|=1$, and this description extends to a finite product by the field K\"unneth formula. Since the algebra is generated by degree one, the description is natural under all linear automorphisms of $V$.

After extending scalars to a splitting field, the eigencharacters of scalar multiplication on $V$ are $\lambda\mapsto\lambda^{2^a}$, $0\leq a<m$, each with multiplicity $\dim_F V$. They may also be read from the isomorphism
$$
F\otimes_{\F_2}F\longrightarrow\prod_{a=0}^{m-1}F,\qquad x\otimes y\longmapsto(x^{2^a}y)_a.
$$
With the contragredient convention the cohomology generators have weights $-2^a$. Dualizing each finite cohomological degree shows that the weights on $H_j(V,\F_2)$ are sums of exactly $j$ members of $\{1,2,4,\ldots,2^{m-1}\}$, with repetitions, modulo $d=2^m-1$.

Quillen's counting lemma \cite[Lemma~16]{Quillen72}, specialized to $p=2$, shows that no such sum is zero modulo $d$ when $0<j<m$. Its cyclic carrying argument is as follows. Place one counter for each summand in the corresponding cyclic position of $\Z/m\Z$. Replace two counters in one position by one counter in the next. Each replacement preserves the residue modulo $d$, including the replacement in the last position, and strictly decreases the positive number of counters. The process terminates at a nonempty set of distinct positions with fewer than $m$ counters. Its ordinary sum lies between $1$ and $d-1$, so its residue is nonzero.

The order $d$ is odd. On a nontrivial character in characteristic two the norm is zero. It follows from the weight calculation that $N_C=0$ on $H_j(V,\F_2)$ for $1\leq j<m$. By the injection~\eqref{eq:integral-mod-two}, the same holds for integral homology. Inverse multiplication negates all the weights and leaves this conclusion unchanged.

For arbitrary $V$, its finite-dimensional $F$-subspaces form a directed system of $C$-stable subgroups. Every finite bar chain is contained in one of them. Since filtered colimits of abelian groups are exact, the bar complex shows that group homology commutes with this union. The exponent and norm assertions therefore hold without a finiteness assumption on $V$.

If $H$ acts $F$-linearly, its action commutes with $C$. The norm acts coefficientwise as zero on the entire bar complex computing $H_i(H,H_j(V,\Z))$. Hence it is zero on its homology, which is annihilated by $2$. Both properties pass to $C$-stable subquotients. On their coinvariants the norm is multiplication by the odd integer $d$, hence the identity. Those coinvariants are therefore zero.
\end{proof}

We will also use exactness of odd-order coinvariants on $2$-primary modules. If $C$ has odd order $d$ and $M$ is a $2$-primary $C$-module, multiplication by $d$ is invertible on $M$. The averaging operator $d^{-1}N_C$ identifies invariants and coinvariants and makes both functors exact on this category. In particular, if $M_C=0$, every $C$-stable subquotient of $M$ has zero coinvariants. For modules annihilated by $2^a$, the inverse of $d$ can be chosen as an integer inverse modulo $2^a$.

Assume now that $A^2\cong A$. Induction on $m$ shows that $A^m\cong A$ for every $m\geq1$: if $A^m\cong A$, then $A^{m+1}\cong A^m\oplus A\cong A\oplus A\cong A$. Represent one such isomorphism by a row $L$ and its inverse by a column $L'$, so that $LL'=1$ and $L'L=I_m$. Then $T\mapsto LTL'$ and $a\mapsto L'aL$ are mutually inverse unital ring isomorphisms between $M_m(A)$ and $A$. Thus $\Gamma_m\cong A^\times$. Composing the regular representation $\F_{2^m}\to M_m(\F_2)$ with the scalar inclusion into $M_m(A)$ and this isomorphism defines a unital embedding $\F_{2^m}\to A$. It is injective because $A\neq0$.

For $k,q\geq1$, put
$$
J_{k,q}(A)=
 \left\{\begin{pmatrix}I_k&B\\0&H\end{pmatrix}:
 B\in M_{k,q}(A)\text{ and }H\in\Gamma_q\right\}. \eqnum[eq:frame-stabilizer]
$$
Write $U=M_{k,q}(A)_{\mathrm{add}}$. The extension $1\to U\to J_{k,q}(A)\xrightarrow{\pi}\Gamma_q\to1$ splits by $s(H)=\operatorname{diag}(I_k,H)$. If $u(B)=\left(\begin{smallmatrix}I_k&B\\0&I_q\end{smallmatrix}\right)$, then $s(H)u(B)s(H)^{-1}=u(BH^{-1})$. Thus the left $\Gamma_q$-action on $U$ is $B\mapsto BH^{-1}$. In particular,
$$
\begin{pmatrix}I_k&B\\0&H\end{pmatrix}
 =u(BH^{-1})s(H)\qquad\text{and}\qquad
 \begin{pmatrix}I_k&B\\0&H\end{pmatrix}^{-1}
 =\begin{pmatrix}I_k&-BH^{-1}\\0&H^{-1}\end{pmatrix}.
$$

Fix $n>0$, choose $m>n$, and embed $F=\F_{2^m}$ into $A$ as above. Let $C=F^\times$. Conjugation by $t_\lambda=\operatorname{diag}(\lambda I_k,I_q)$ fixes $s(\Gamma_q)$ pointwise and sends $B\in U$ to $\lambda B$. This scalar action commutes with the $\Gamma_q$-action because $\lambda(BH^{-1})=(\lambda B)H^{-1}$.

The next calculation is a form of the scalar argument for affine groups; compare \cite[Theorem~1.9]{Suslin84}, \cite[Proposition~1.10, Theorem~1.11 and Remark~1.13]{NS90}, \cite[\S2]{Schlichting17}, and \cite[Proposition~2.1]{Knudson02}. The role of central scalars in the Nesterenko--Suslin argument is noted in \cite[Definition~2.1, footnote~1]{Schlichting17}. We do not assume that $F$ is central in $A$. Left scalar multiplication commutes with the action $B\mapsto BH^{-1}$, as required in the proof.

\begin{lemma}\label{lem:stabilizer-scalars}
For the scalar action just defined, let
$$
K_n=\ker\bigl(H_n(J_{k,q}(A),\Z)\xrightarrow{\pi_*}H_n(\Gamma_q,\Z)\bigr).
$$
Then $2^nK_n=0$, $(K_n)_C=0$, and the projection induces an isomorphism
$$
H_n(J_{k,q}(A),\Z)_C\cong H_n(\Gamma_q,\Z). \eqnum[eq:stabilizer-coinvariants]
$$
If the standard inclusion $\Gamma_q\to\Gamma_{k+q}$, $H\mapsto\operatorname{diag}(H,I_k)$, induces zero on $H_n(-,\Z)$, then the inclusion $J_{k,q}(A)\to\Gamma_{k+q}$ also induces zero on $H_n(-,\Z)$. All these assertions remain valid when the scalar actions are inverted.
\end{lemma}

\begin{proof}
Associated with the split extension is the $C$-equivariant Lyndon--Hochschild--Serre spectral sequence
$$
E^2_{p,j}=H_p(\Gamma_q,H_j(U,\Z)) \ \Longrightarrow\ H_{p+j}(J_{k,q}(A),\Z). \eqnum[eq:stabilizer-lhs]
$$
We use its usual homological indexing and projection edge map; see \cite[Theorem~6.8.2]{Weibel94}. The base row is permanent. Indeed, the section and projection compare it with the spectral sequence of the extension of $\Gamma_q$ with trivial kernel and induce the identity on that row. There are no incoming differentials, and naturality with the section makes the outgoing differentials zero.

The kernel $K_n$ is the penultimate filtration subgroup of $H_n(J_{k,q}(A),\Z)$. Its graded pieces are $E^\infty_{p,n-p}$ for $0\leq p\leq n-1$. Each is a $C$-stable subquotient of a term in~\eqref{eq:stabilizer-lhs} with $1\leq n-p\leq n<m$. Lemma~\ref{lem:scalar-homology} shows that these pieces are annihilated by $2$ and have zero $C$-coinvariants. There are at most $n$ pieces, so $2^nK_n=0$. Right exactness of coinvariants, applied successively to the filtration, implies $(K_n)_C=0$. This argument includes all mixed differentials, since every limiting term is a subquotient of its $E^2$ term.

The section splits $\pi_*$ and is fixed by $C$, proving~\eqref{eq:stabilizer-coinvariants}. Let $i:J_{k,q}(A)\to\Gamma_{k+q}$ be the inclusion. Its composite with a scalar automorphism is its composite with ambient conjugation by $t_\lambda$. Inner automorphisms induce the identity on homology, so $i_*$ factors through~\eqref{eq:stabilizer-coinvariants}. Its image is therefore the image of the inclusion $H\mapsto\operatorname{diag}(I_k,H)$. A block permutation conjugates this map to $H\mapsto\operatorname{diag}(H,I_k)$, which is the specified standard inclusion. Its map on $H_n$ is zero by hypothesis. Inverting the scalar actions permutes the summands of the norm and preserves the argument.
\end{proof}

\begin{proof}[Proof of Theorem~\ref{thm:acyclicity-criterion}]
The group $\Gamma_r$ acts transitively on the ordered $(k-1)$-simplices of $X_r(A)$. An isomorphism $A^{r-k}\to C$ in~\eqref{eq:frame-decomposition} completes a given frame to an invertible matrix. The stabilizer of the standard $k$-frame is $J_{k,r-k}(A)$, and the stabilizer of a full frame is trivial. We put $J_{r,0}(A)=1$. Since the frames are ordered, the stabilizers fix their simplices pointwise.

Let $\mathcal B_r=E\Gamma_r\times_{\Gamma_r}|X_r(A)|$. For a free right $\Z\Gamma_r$-resolution $P_*$ of $\Z$, the double complex $P_t\otimes_{\Z\Gamma_r}C_p(X_r(A),\Z)$ computes $H_*(\mathcal B_r,\Z)$. Its two filtrations induce the Borel homology comparison and the stabilizer spectral sequence
$$
E^1_{p,t}=H_t(J_{p+1,r-p-1}(A),\Z) \ \Longrightarrow\ H_{p+t}(\mathcal B_r,\Z), \qquad\text{for }0\leq p\leq r-1; \eqnum[eq:frame-spectral-sequence]
$$
see \cite[\S6.1.15]{Weibel94}. This is the standard construction from frame stabilizers; compare \cite[Lemmas~2.3--2.4]{NS90}. All filtrations in a fixed total degree are finite. By the second hypothesis, the projection $\mathcal B_r\to B\Gamma_r$ induces an isomorphism on $H_n$ when $r\geq n+3$: in the other spectral sequence, $H_a(\Gamma_r,H_b(X_r(A),\Z))$, all positive fiber rows through degree $n$ vanish, and the row in degree zero is ordinary group homology.

The row $t=0$ in~\eqref{eq:frame-spectral-sequence} has one copy of $\Z$ in each column. The differential out of column $p$ is multiplication by $\sum_{i=0}^p(-1)^i$. Thus
$$
E^2_{p,0}=0\qquad\text{for }1\leq p\leq r-2. \eqnum[eq:frame-bottom-row]
$$
For positive even $p$ the outgoing differential is injective, and for odd $p$ the incoming differential from $p+1$ is the identity.

We induct on $n$. Assume $H_t(A^\times,\Z)=0$ for $0<t<n$; for $n=1$ this assumption is empty. Choose $r=n+3$, choose $m>n$, and fix one unital copy of $F=\F_{2^m}$ in $A$. The group $C=F^\times$ acts on the whole frame semisimplicial set by
$$
T_\lambda(v_1,\ldots,v_k)=(v_1\lambda,\ldots,v_k\lambda). \eqnum[eq:global-rescaling]
$$
Each $\lambda$ is a two-sided unit, so the summands and witness complements in~\eqref{eq:frame-decomposition} are unchanged. The maps commute with deletion and are $\Gamma_r$-equivariant, since $g(v_i\lambda)=(gv_i)\lambda$. The composition law holds because $C$ is abelian.

The identity on $E\Gamma_r$ and~\eqref{eq:global-rescaling} on $X_r(A)$ define an action on $\mathcal B_r$ covering the identity of $B\Gamma_r$. On the double complex it is the strict chain action $1\otimes T_\lambda$, commuting with both differentials and preserving the filtration. Since the projection is an isomorphism on $H_n$, this action is trivial on $H_n(\mathcal B_r,\Z)$.

To determine the action on~\eqref{eq:frame-spectral-sequence}, let $\sigma_k$ be the standard $k$-frame and put $t_{\lambda,k}=\operatorname{diag}(\lambda I_k,I_{r-k})$. Then $T_\lambda\sigma_k=t_{\lambda,k}\sigma_k$. Transport back to $\sigma_k$ by $t_{\lambda,k}^{-1}$. The resulting stabilizer automorphism is
$$
h\longmapsto t_{\lambda,k}^{-1}ht_{\lambda,k},\qquad\text{that is,}\qquad
 \begin{pmatrix}I_k&B\\0&H\end{pmatrix}
 \longmapsto\begin{pmatrix}I_k&\lambda^{-1}B\\0&H\end{pmatrix}.
$$
Equivalently, the orbit map $gJ\mapsto gt_{\lambda,k}J$ induces the map $x\mapsto xt_{\lambda,k}$ on $E\Gamma_r/J$, and $(xh)t_{\lambda,k}=(xt_{\lambda,k})(t_{\lambda,k}^{-1}ht_{\lambda,k})$ confirms the conjugation convention. Thus the $E^1$ action is the inverse scalar action of Lemma~\ref{lem:stabilizer-scalars}.

The chosen deletion representatives need only commute with these stabilizer automorphisms up to inner conjugation. If the rows of $B$ are $b_1,\ldots,b_k$, deleting row $i$ and placing that coordinate first among the new quotient coordinates defines the map
$$
f_i\begin{pmatrix}I_k&B\\0&H\end{pmatrix}
 =\begin{pmatrix}
 I_{k-1}&0&B_{\widehat i}\\
 0&1&b_i\\
 0&0&H
 \end{pmatrix}.
$$
The source scalar action multiplies both $B_{\widehat i}$ and $b_i$ by $\lambda^{-1}$; the target scalar action multiplies only $B_{\widehat i}$. Their discrepancy is conjugation by $\operatorname{diag}(I_{k-1},\lambda^{-1},I_q)$ in the target stabilizer, where $q=r-k$. The strict global chain action above supplies equivariance on every page of the spectral sequence, including this coordinate comparison.

For $0<t<n$, the isomorphism $\Gamma_q\cong A^\times$ and the induction hypothesis imply $H_t(\Gamma_q,\Z)=0$. By Lemma~\ref{lem:stabilizer-scalars}, $H_t(J_{k,q}(A),\Z)$ is annihilated by $2^t$ and has zero $C$-coinvariants. This also holds for the trivial stabilizer with $q=0$. Since coinvariants under a group of odd order are exact on $2$-primary modules,
$$
\bigl(E^\infty_{p,n-p}\bigr)_C=0\qquad\text{for }1\leq p<n. \eqnum[eq:lower-filtration-coinvariants]
$$
Every limiting term here is a $C$-stable subquotient of its $E^1$ term, irrespective of the differentials reaching it.

The term $E^\infty_{n,0}$ is zero by~\eqref{eq:frame-bottom-row}. The remaining vertex term is also zero. Indeed, the filtration subgroup $F_0H_n(\mathcal B_r,\Z)$ is the image of $H_n(J_{1,r-1}(A),\Z)$. Its composite with the projection to $H_n(\Gamma_r,\Z)$ is the map induced by the stabilizer inclusion. This map is zero by the first hypothesis and Lemma~\ref{lem:stabilizer-scalars}, with $k=1$ and $q=n+2$. Since that projection is an isomorphism, $F_0H_n(\mathcal B_r,\Z)=0$. This calculation uses the limiting image, with all differentials into the vertex term already accounted for.

All graded pieces of the finite filtration of $H_n(\mathcal B_r,\Z)$ have zero $C$-coinvariants, by~\eqref{eq:lower-filtration-coinvariants} and the two endpoint calculations. By right exactness, applied successively to this filtration, $H_n(\mathcal B_r,\Z)_C=0$. The action on this group is trivial, so $H_n(\mathcal B_r,\Z)=0$. The Borel projection and $\Gamma_r\cong A^\times$ now imply $H_n(A^\times,\Z)=0$, completing the induction.
\end{proof}

\section{Leaf coordinates and standard inclusions}
\label{sec:leaf-coordinates}

We apply the criterion to the binary Leavitt algebra
$$
R=\F_2\langle e,f,e^*,f^*\rangle/ (e^*e=f^*f=1,\ e^*f=f^*e=0,\ ee^*+ff^*=1), \eqnum[eq:leavitt-algebra]
$$
and write $G_r=\GL_r(R)$ and $G=R^\times$. A positive word is a word in $e,f$, including the empty word. If $\alpha=c_1\cdots c_h$, write $\alpha^*=c_h^*\cdots c_1^*$.

The algebra $R$ is nonzero. On the $\F_2$-vector space with basis the infinite binary words, let $e,f$ prepend the indicated letter. Let $e^*,f^*$ remove a matching first letter and send a word beginning with the other letter to zero. These operators satisfy~\eqref{eq:leavitt-algebra}, with nonzero identity.

For these standard coordinates in Leavitt algebras, see \cite{AAP08,AAS17}. We shall use the following inverse maps.

A complete finite binary leaf set is obtained from the empty word by repeatedly replacing a word $\alpha$ by $\alpha e,\alpha f$. The order of its leaves may be chosen arbitrarily.

\begin{lemma}\label{lem:leaf-coordinates}
For a complete ordered leaf set $(\alpha_1,\ldots,\alpha_m)$, let $L=(\alpha_1,\ldots,\alpha_m)$ and $L^*=(\alpha_1^*,\ldots,\alpha_m^*)^{\mathrm t}$. Then
$$
LL^*=1\qquad\text{and}\qquad L^*L=I_m. \eqnum[eq:leaf-identities]
$$
Consequently $x\mapsto Lx$ and $a\mapsto L^*a$ are mutually inverse maps of right modules $R^m\rightleftarrows R$, and
$$
M_m(R)\rightleftarrows R,\qquad Z\longmapsto LZL^*\qquad\text{and}\qquad a\longmapsto L^*aL \eqnum[eq:leaf-ring-isomorphisms]
$$
are mutually inverse unital ring isomorphisms. Every positive integer $m$ occurs as the size of a complete leaf set.
\end{lemma}

\begin{proof}
Distinct leaves are prefix-incomparable. Cancellation therefore implies $\alpha_i^*\alpha_j=\delta_{ij}$. Starting with $1=1\cdot1^*$, each leaf replacement preserves the sum of the range projections, because $\alpha\alpha^*=\alpha ee^*\alpha^*+\alpha ff^*\alpha^*$. Hence $\sum_i\alpha_i\alpha_i^*=1$, proving~\eqref{eq:leaf-identities}. These identities verify the two pairs of inverse maps and show that the ring maps preserve products and units. A complete leaf set of size $m$ is the comb
$$
L_m=(e^{m-1},e^{m-2}f,\ldots,ef,f),\qquad\text{with }L_1=(1).
$$
Its recurrence is $L_{m+1}=(eL_m,f)$.
\end{proof}

To prove that the standard inclusions induce zero on homology, we combine the stable calculation of \cite{ABC09} with a finite compression identity. Sum rings occur in algebraic $K$-theory \cite[\S2]{Wagoner72}; here the compression uses only finite sums.

\begin{proposition}\label{prop:zero-padding}
For every $r\geq1$, the standard inclusion $G_r\to G_{r+1}$, $g\mapsto\operatorname{diag}(g,1)$, induces zero on $H_n(-,\Z)$ for every $n>0$.
\end{proposition}

\begin{proof}
For the graph with one vertex and two loops, the algebraic $K$-theory sequence of Ara--Brustenga--Corti\~nas \cite[Theorem~7.6]{ABC09} has map $1-2=-1$ on $K(\F_2)$. The graph is row-finite and has no sinks. The coefficient field is regular supercoherent: every polynomial ring in finitely many variables over $\F_2$ is regular Noetherian. By the theorem, $K(R)$ is the homotopy cofiber of an equivalence. Hence $K_i(R)=0$ for all $i>0$.

Let $\GL_\infty(R)$ be the direct limit under standard inclusions. Its plus construction has fundamental group $K_1(R)$ and higher homotopy groups $K_i(R)$, $i\geq2$, and preserves integral homology \cite[Chapter~IV, Definitions~1.1 and~1.1.1]{Weibel13}. It is a connected CW complex with trivial homotopy groups, hence is contractible by Whitehead's theorem. Thus
$$
H_n(\GL_\infty(R),\Z)=0\qquad\text{for }n>0. \eqnum[eq:stable-homology-zero]
$$

Define $c:G\to G$ by $c(u)=eue^*+ff^*$. By orthogonality, $c(u)c(v)=c(uv)$, $c(1)=1$, and $c(u)^{-1}=c(u^{-1})$. For $s\geq0$, let $\Phi_s:R\to M_{2^s}(R)$ have entries $\Phi_s(a)_{\alpha,\beta}=\alpha^*a\beta$, where $\alpha,\beta$ range over the words of length $s$. These are the ring isomorphisms of Lemma~\ref{lem:leaf-coordinates}. Order the words at depth $s+1$ by $e\alpha$ followed by $f\alpha$. Multiplying the coordinates, we obtain
$$
\Phi_{s+1}(c(u))=\operatorname{diag}(\Phi_s(u),I_{2^s}). \eqnum[eq:dyadic-compression]
$$

Thus the system $(G,c)$ is isomorphic to the dyadic cofinal subsystem of the standard general linear system. The normalized bar complex commutes with this filtered direct limit, since every generator and every finite chain occurs at a finite stage. According to~\eqref{eq:stable-homology-zero}, $\varinjlim(H_n(G,\Z),c_*)=0$. Therefore each element of $H_n(G,\Z)$ is killed by a positive iterate of $c_*$.

The units associated with the complete leaf sets $(e,fe,ff)$ and $(ee,ef,f)$ are
$$
\begin{aligned}
 w&=ee\,e^*+ef(fe)^*+f(ff)^*\qquad\text{and}\\
 w^{-1}&=e(ee)^*+fe(ef)^*+ff\,f^*.
 \end{aligned}
$$
Both inverse products follow from~\eqref{eq:leaf-identities}. Moreover $we=ee$ and $w(ff^*)w^{-1}=ef(ef)^*+ff^*$. Hence
$$
wc(u)w^{-1}=ee\,u(ee)^*+ef(ef)^*+ff^*=c(c(u)). \eqnum[eq:compression-idempotence]
$$
Inner conjugation acts trivially on homology, so $c_*^2=c_*$. This idempotent is locally nilpotent by the preceding paragraph, and consequently $c_*=0$ on every positive homology group.

For $Z\in G_r$, the comb recurrence takes the form
$$
L_{r+1}\operatorname{diag}(Z,1)L_{r+1}^* =c(L_rZL_r^*). \eqnum[eq:padding-compression]
$$
The two leaf coordinate maps are group isomorphisms. Since $c_*=0$, this identity proves the proposition.
\end{proof}

The map $c$ used here is a group homomorphism. Its affine formula sends $0$ to $ff^*\neq0$, so it is not a ring homomorphism. Proposition~\ref{prop:zero-padding} concerns the specified standard inclusions; the coordinate isomorphisms of Lemma~\ref{lem:leaf-coordinates} induce homology isomorphisms.

\section{Simultaneous extensions of ordered frames}
\label{sec:common-cones}

It remains to verify the homology hypothesis in Theorem~\ref{thm:acyclicity-criterion}. We use a common transverse vector to cone each finite family of frames; see \cite[\S2.7, Remark~4]{vdK80} and \cite[\S2, Lemmas~2.1--2.2]{Suslin84}. For $R$, the required extension follows from a simultaneous version of Leavitt word reduction. The reduction theorem for one nonzero element is proved in \cite[Lemma~2.2.8, Theorem~2.2.11 and Corollary~2.2.12]{AAS17}. We need a single right multiplier for a finite family, with nonzero kernels and specified finite bases for the resulting left inverses.

If $\mu_1,\ldots,\mu_s$ are distinct positive words of lengths at most $D$ and $M>D$, then the words $\mu_i e^M f$ are pairwise prefix-incomparable. Indeed, for comparable original words write $\mu_j=\mu_i\delta$ with $\delta$ nonempty. If $\delta$ contains $f$, its first $f$ occurs before position $M+1$; if $\delta=e^a$, the first $f$ in $\delta e^M f$ occurs at position $M+a+1$. In either case $e^M f$ and $\delta e^M f$ disagree before the shorter word ends. Original words that were prefix-incomparable remain so. Therefore
$$
(\mu_i e^M f)^*(\mu_j e^M f)=\delta_{ij}. \eqnum[eq:word-separation]
$$
In particular, a nonempty sum of distinct positive words is nonzero in $R$: append $e^M f$ and multiply on the left by the star of one of the resulting words to obtain $1$.

\begin{lemma}\label{lem:word-multiplier}
Given finitely many nonzero $a_1,\ldots,a_s\in R$, there are $x\in R$ and nonempty positive words $\eta_1,\ldots,\eta_s$ such that $\eta_i^*a_ix=1$ for every $i$. Each decomposition $R=a_ixR\oplus\ker\eta_i^*$ has an explicit nonzero complement with a finite basis consisting of positive words.
\end{lemma}

\begin{proof}
Every element of $R$ is a finite sum of monomials $\alpha\beta^*$. To obtain this form, cancel every occurrence of a starred letter followed by an unstarred letter. Choose expressions of this form for the $a_i$, and choose $N$ at least as large as all lengths $|\beta|$ occurring in these expressions. Let $\gamma_1,\ldots,\gamma_t$ be all words of length $N$. Each $a_i\gamma_j$ is a polynomial in positive words. For every $i$, at least one is nonzero, since $a_i=\sum_j(a_i\gamma_j)\gamma_j^*$. Collect identical words over $\F_2$ in all these polynomials, and let $L$ bound the lengths of the remaining words.

Put $t_j=e^{j(L+1)}f$ and $x_0=\sum_j\gamma_jt_j$. The words contributed by $(a_i\gamma_j)t_j$ have lengths in $[j(L+1)+1,j(L+1)+1+L]$. These intervals are disjoint for distinct $j$, and concatenation by a fixed $t_j$ cannot identify distinct words within one interval. Thus $P_i=a_ix_0$ is a nonempty positive polynomial for each $i$.

Choose $M$ greater than every word length occurring in the $P_i$ and set $x=x_0e^Mf$. Equation~\eqref{eq:word-separation} makes the support of each $P_ie^Mf$ a prefix antichain. Choose $\eta_i$ in this support. Its coefficient is $1$, so $\eta_i^*a_ix=1$ by~\eqref{eq:word-separation}. Every $\eta_i$ is nonempty.

For a word $\eta=c_1\cdots c_h$, $h\geq1$, let $s_\ell=c_1\cdots c_{\ell-1}c'_\ell$, where $c'_\ell$ is the other binary letter. The words $\eta,s_1,\ldots,s_h$ form a complete leaf set. Consequently
$$
\ker\eta^*=\bigoplus_{\ell=1}^h s_\ell R, \eqnum[eq:sibling-kernel]
$$
with inverse coordinate maps $(r_\ell)_\ell\mapsto\sum_\ell s_\ell r_\ell$ and $z\mapsto(s_\ell^*z)_\ell$. If $y=a_ix$ and $b=\eta_i^*$, then $by=1$, and
$$
R\oplus\ker b\rightleftarrows R,\qquad (t,z)\longmapsto yt+z\qquad\text{and}\qquad u\longmapsto(bu,u-ybu)
$$
are mutually inverse. The kernel in~\eqref{eq:sibling-kernel} is nonzero because $h\geq1$ and $s_\ell^*s_\ell=1$.
\end{proof}

\begin{proposition}\label{prop:frame-cone}
For a finite family of simplices of $X_r(R)$, each having at most $r-2$ frame vectors, there is a single vector $v$ extending every frame in the family to a simplex of $X_r(R)$. Each extended frame has a nonzero complement with explicit coordinates of the prescribed formal size.
\end{proposition}

\begin{proof}
For an empty family, choose any standard vertex. Otherwise, let $W_i$ be the span of the $i$th frame, with its given ordered basis, let $k_i$ be its length, and put $q_i=r-k_i\geq2$. Choose a completion $R^r=W_i\oplus C_i^0$ with $C_i^0\cong R^{q_i}$. Compose the quotient coordinate map with $L_{q_i}:R^{q_i}\to R$ to obtain $\rho_i:R^r\to R$ with kernel $W_i$. Let $\sigma_i:R\to R^r$ be the corresponding section with image $C_i^0$, so that $\rho_i\sigma_i=1$.

Fix $\Phi=L_r^*:R\to R^r$. The endomorphism $\rho_i\Phi$ of the right module $R$ is left multiplication by an element $a_i\neq0$, because $\rho_i$ is surjective and $\Phi$ is an isomorphism. Apply Lemma~\ref{lem:word-multiplier} to the $a_i$, and put $v=\Phi x$, $y_i=\rho_i v=a_ix$, $b_i=\eta_i^*$, and $w_i=v-\sigma_i y_i\in W_i$. Since $b_i\rho_i v=1$, the map $R\to R^r$, $a\mapsto va$, is a split injection.

There is a decomposition
$$
R^r=W_i\oplus vR\oplus\sigma_i\ker b_i. \eqnum[eq:cone-decomposition]
$$
Its coordinate map and inverse are
$$
\begin{aligned}
 (w,t,z)&\longmapsto w+vt+\sigma_i z\qquad\text{and}\\
 u&\longmapsto\begin{pmatrix}
 (I-\sigma_i\rho_i)u-w_i b_i\rho_i u\\
 b_i\rho_i u\\
 \rho_i u-y_i b_i\rho_i u
 \end{pmatrix},
 \end{aligned} \eqnum[eq:cone-coordinate-inverses]
$$
where $w\in W_i$, $t\in R$, and $z\in\ker b_i$. The first inverse component belongs to $W_i$ and the last to $\ker b_i$. The image of the forward map under $\rho_i$ is $y_it+z$, so the second and third inverse components are $t$ and $z$; the first is then $w$. Conversely, substitution of the three inverse components into the forward map returns $u$, since $v=\sigma_i y_i+w_i$. This verifies both inverse products.

We now identify the complement with the free module of the required formal size. Let $h_i=|\eta_i|$, and let $S_i=(s_1,\ldots,s_{h_i})$ be the sibling row in~\eqref{eq:sibling-kernel}. The map
$$
\sigma_iS_iL_{h_i}^*L_{q_i-1}:R^{q_i-1} \longrightarrow\sigma_i\ker b_i \eqnum[eq:complement-formal-coordinates]
$$
is an isomorphism, with inverse on this complement $L_{q_i-1}^*L_{h_i}S_i^*\rho_i$. By $\rho_i\sigma_i=1$, $S_i^*S_i=I_{h_i}$, and~\eqref{eq:leaf-identities}, the inverse product on $R^{q_i-1}$ is the identity. The other product is $\sigma_iS_iS_i^*\rho_i$, which is the identity on $\sigma_i\ker b_i$ by~\eqref{eq:sibling-kernel}. Since $q_i-1\geq1$ and $h_i\geq1$, the complement is nonzero. Thus~\eqref{eq:cone-decomposition} extends every specified frame at the same formal rank $r$.
\end{proof}

\begin{corollary}\label{cor:frame-homology}
For $r\geq3$, one has $\widetilde H_d(X_r(R),\Z)=0$ for $0\leq d\leq r-3$.
\end{corollary}

\begin{proof}
Let $z$ be a finite ordered integral $d$-cycle, using the reduced chain complex when $d=0$. Each frame in its support has $d+1\leq r-2$ vectors. Choose a common extension $v$ by Proposition~\ref{prop:frame-cone}. Prepending $v$ to every support simplex, and to its faces, defines a chain operator $h_v$. For the ordered boundary, $\partial h_v+h_v\partial=\id$ on these chains. At the augmentation, set $h_v(1)=[v]$. All signs are the usual integral alternating signs. Since $\partial z=0$, one has $\partial h_vz=z$. The same argument for reduced zero-cycles also proves connectedness.
\end{proof}

\begin{theorem}\label{thm:acyclicity}
The unit group of $R=L_{\F_2}(1,2)$ is integrally acyclic. More generally, $H_n(\GL_r(R),\Z)=0$ for every $n>0$ and $r\geq1$.
\end{theorem}

\begin{proof}
Lemma~\ref{lem:leaf-coordinates}, Proposition~\ref{prop:zero-padding}, and Corollary~\ref{cor:frame-homology} verify the hypotheses of Theorem~\ref{thm:acyclicity-criterion}. It follows that $H_n(R^\times,\Z)=0$ for all $n>0$. The leaf ring isomorphisms~\eqref{eq:leaf-ring-isomorphisms} identify every $\GL_r(R)$ with $R^\times$, so the same conclusion holds in every rank.
\end{proof}

\section{The Steinberg comparison}
\label{sec:steinberg}

The acyclicity of the unit group will control the kernel of the Steinberg map. We first relate this kernel to the topology of the ordered frame complex over a ring of characteristic two.

For a nonzero unital ring $B$, write $t_{ij}(a)=I+aE_{ij}$ and let $E_m(B)$ be the subgroup of $\GL_m(B)$ generated by these matrices. For $m\geq3$, the group $\St_m(B)$ has generators $X_{ij}(a)$, where $i\ne j$ and $a\in B$, and relations
$$
\begin{aligned}
\eqrownum &\quad X_{ij}(a)X_{ij}(b)=X_{ij}(a+b),\\
\eqrownum &\quad [X_{ij}(a),X_{kl}(b)]=1 \quad \text{for }i\ne l\text{ and }j\ne k,\\
\eqrownum &\quad [X_{ij}(a),X_{jk}(b)]=X_{ik}(ab) \quad \text{for distinct }i,j,k.
\end{aligned}
$$
Here $[x,y]=xyx^{-1}y^{-1}$. Let $\phi_m:\St_m(B)\to\GL_m(B)$ send $X_{ij}(a)$ to $t_{ij}(a)$, and put $N_m(B)=\ker\phi_m$. The standard stabilization homomorphism $j_m:\St_m(B)\to\St_{m+1}(B)$ preserves every root label and coefficient.

In defining the ordered semisimplicial set $X_n(B)$, we keep the formal rank fixed. A $(k-1)$-simplex, for $1\leq k\leq n$, is an ordered tuple $(v_1,\ldots,v_k)$ admitting a decomposition
$$
B^n=v_1B\oplus\cdots\oplus v_kB\oplus C,\qquad\text{with }C\cong B^{n-k},
$$
in which $b\mapsto v_i b$ is an isomorphism $B\to v_iB$. The complement and its coordinates are not part of the simplex. Faces delete a vector. The deleted summand, with its specified basis vector, together with $C$ has the required formal type for the new complement. Thus the face maps are well defined without cancellation of free ranks. There are no simplices with $k>n$.

\begin{theorem}[Steinberg comparison criterion]
\label{thm:raw-criterion}
Let $B$ be a nonzero unital ring of characteristic two, and fix $n\geq4$. Suppose that
\begin{enumerate}
\item $\GL_{n-1}(B)=E_{n-1}(B)$ and $\GL_{n-2}(B)=E_{n-2}(B)$;
\item $j_{n-1}(N_{n-1}(B))=1$ in $\St_n(B)$;
\item $|X_n(B)|$ is simply connected.
\end{enumerate}
Then $\phi_n:\St_n(B)\to\GL_n(B)$ is an isomorphism.
\end{theorem}

\begin{proof}
Put $G=\GL_n(B)$, $S=\St_n(B)$, and let $b_1,\ldots,b_n$ be the standard columns. The action of $G$ on ordered $k$-frames is transitive, since a basis of a permitted complement completes the frame to an automorphism of $B^n$. In particular, there is one orbit in each of dimensions zero, one, and two.

The vertex stabilizer $J$ and the stabilizer $K$ of the ordered edge are
$$
J=\left\{j(b,H)=\begin{pmatrix}1&b\\0&H\end{pmatrix}\right\}\qquad\text{and}\qquad
 K=\left\{k(a,b,H)=
 \begin{pmatrix}1&0&a\\0&1&b\\0&0&H\end{pmatrix}\right\}.
$$
In $J$, the row $b$ has length $n-1$ and $H\in\GL_{n-1}(B)$. In $K$, the rows have length $n-2$ and $H\in\GL_{n-2}(B)$. The upper rows in these full stabilizers are unrestricted. Let $\tau=(12)$ and $h=(23)$ be permutation matrices, so that $h\in J$, and put $\eta(k)=\tau k\tau^{-1}$. This automorphism of $K$ interchanges its two upper rows.

We apply Brown's presentation for a group acting on a simply connected complex \cite[Theorem~1 and \S3]{Brown84}. For the ordered frame complex, the Borel construction $Y=EG\times_G|X_n(B)|$ has fundamental group
$$
\Pi=\langle J,T\mid TkT^{-1}=\eta(k)\text{ for }k\in K,\quad ThT=hTh\rangle. \eqnum[eq:borel-presentation]
$$
This notation includes every multiplication relation in $J$. To verify that the presentation is complete, form the cellular Borel construction through total dimension two. The vertex orbit contributes $BJ$. The edge orbit contributes a cylinder over $BK$: the cell of dimension zero in the base accounts for $T$, and the cells of dimension one in the base impose the conjugation relations between the two endpoint inclusions. The triangle orbit contributes one further cell of dimension two, arising from the cell of dimension zero in the classifying space of its stabilizer. All other cells arising from triangle stabilizers, and all cells arising from simplices of dimension at least three, have total dimension at least three. The cells of dimension two associated with the vertex stabilizer impose precisely the relations already included in $J$.

To determine the attaching relation with its orientation, identify the fundamental group of $Y$ with the group of pairs $(g,[p])$, where $p$ is a path from $b_1$ to $gb_1$, taken up to homotopy relative to its endpoints. Multiplication is $(g,[p])(g',[p'])=(gg',[p*(gp')])$. Represent $j\in J$ by its constant path and $T$ by $(\tau,E)$, where $E$ is the ordered edge $(b_1,b_2)$. For $k\in K$, both $Tk$ and $\eta(k)T$ have projection $\tau k$ and path $E$. The word $ThT$ has projection $(13)$ and path $b_1\to b_2\to b_3$, whereas $hTh$ has the same projection and path $b_1\to b_3$. The boundary of the ordered triangle $(b_1,b_2,b_3)$ therefore imposes the second relation in \eqref{eq:borel-presentation}.

The same path description also verifies completeness directly. Every edge traversal is a translate of $E$ or its reverse. Starting at a vertex represented by a transporter $g$, it is encoded by a word $jT$ or $jT^{-1}$, with $j\in J$. Two choices describing the same ordered edge differ by its stabilizer $K$, and the first relations in \eqref{eq:borel-presentation} identify them. At the endpoint, the difference from a prescribed transporter is again in $J$. Homotopies of edge paths are generated by backtracking and passages across individual triangles. Backtracking cancels inverse words, and a passage across a triangle is a translate of the relation just computed. Thus no additional path relation is required. In particular, the two ordered edge cells $(u,v)$ and $(v,u)$ remain distinct. The relation $T^2=1$ follows from the displayed presentation: since $h^2=1$, the braid relation implies $(hT)h(hT)^{-1}=T$.

The homotopy sequence of $|X_n(B)|\to Y\to BG$, together with the third hypothesis, shows that the projection $p:\Pi\to G$, given by $p(j)=j$ and $p(T)=\tau$, is an isomorphism.

We lift the relations of this presentation to $S$. Permuting indices is an automorphism of the Steinberg presentation, so the second hypothesis also applies to stabilization on coordinates $2,\ldots,n$. By the first hypothesis, we can define a homomorphism $\ell:\GL_{n-1}(B)\to S$ covering $H\mapsto\operatorname{diag}(1,H)$ by choosing an elementary word for $H$ and evaluating it on these coordinates. Two choices differ by a kernel word killed by the second hypothesis, so the homomorphism is well defined.

For a row $b=(b_2,\ldots,b_n)$ set $x_1(b)=\prod_{j=2}^nX_{1j}(b_j)$, in increasing order. The factors commute. According to the Steinberg relations,
$$
\ell(H)x_1(b)\ell(H)^{-1}=x_1(bH^{-1}). \eqnum[eq:row-action]
$$
Indeed, for $H=t_{ij}(c)$ with $i,j\geq2$, conjugation changes only the $j$-th entry, replacing it by $b_j-b_i c$. This is exactly right multiplication by $I-cE_{ij}$; this term comes from the third Steinberg relation, while all other factors commute. Elementary generation extends the identity to every $H$.

Define $\sigma(j(b,H))=\ell(H)x_1(b)$. The chosen order of the factors preserves the upper row $b$. It follows from \eqref{eq:row-action} that
$$
\sigma(j(b,H))\sigma(j(b',H')) =\ell(HH')x_1(bH'+b') =\sigma(j(bH'+b',HH')),
$$
which is the matrix multiplication law in $J$. Hence $\sigma:J\to S$ is a homomorphism lifting its inclusion in $G$.

For $H\in\GL_{n-2}(B)$, elementary generation allows us to choose a lift $\ell_0(H)$ supported on coordinates $3,\ldots,n$. It is independent of the elementary word, since a null word in these coordinates is a kernel word in a block of rank $(n-1)$ and is then killed by the stabilization hypothesis. For $n=4$, this comparison takes place in $\St_3(B)$ on coordinates $2,3,4$. In particular $\ell_0(H)=\ell(\operatorname{diag}(1,H))$. For the element $k(a,b,H)\in K$ one consequently has
$$
\sigma(k(a,b,H))=\ell_0(H)x_2(b)x_1(a), \eqnum[eq:edge-section]
$$
where both row products now range over columns $3,\ldots,n$. The second equality in the first hypothesis is used here to lift the subgroup $\operatorname{diag}(I_2,\GL_{n-2}(B))$.

Put $w_{ij}=X_{ij}(1)X_{ji}(1)X_{ij}(1)$. By additivity in characteristic two, $w_{ij}^2=1$. The Steinberg relations also imply
$$
w_{12}X_{ij}(c)w_{12}^{-1} =X_{\tau(i),\tau(j)}(c). \eqnum[eq:weyl-action]
$$
We verify this identity for all generators. Write $u=X_{12}(1)$ and $v=X_{21}(1)$. For $j\geq3$, the successive conjugates of $X_{1j}(c)$ by $u,v,u$ are $X_{1j}(c)$, then $X_{2j}(c)X_{1j}(c)$, and finally $X_{2j}(c)$, since the two copies of $X_{1j}(c)$ cancel. Starting with $X_{j1}(c)$, the successive conjugates are $X_{j2}(c)X_{j1}(c)$, then $X_{j2}(c)$, and again $X_{j2}(c)$. The reverse identities follow from $w_{12}^2=1$. Generators $X_{ij}(c)$ with $i,j\notin\{1,2\}$ commute with $u$ and $v$. Finally, express $X_{12}(c)$ as $[X_{1j}(c),X_{j2}(1)]$, with $j\geq3$, and use the identities just proved; the result is $X_{21}(c)$. This also treats $X_{21}(c)$ and proves \eqref{eq:weyl-action}.

It follows that $w_{12}$ commutes with $\ell_0(H)$ and interchanges the two row products in \eqref{eq:edge-section}. Those products commute with each other. Therefore $w_{12}\sigma(k)w_{12}^{-1}=\sigma(\eta(k))$ for every $k\in K$. Moreover $\sigma(h)=w_{23}$, because the indicated elementary word on the last coordinates represents the permutation $(23)$. On applying \eqref{eq:weyl-action} to the three factors of a Weyl word, we obtain $w_{12}w_{23}w_{12}=w_{13}=w_{23}w_{12}w_{23}$. Thus the triangle relation also holds in $S$.

The presentation \eqref{eq:borel-presentation} therefore defines a homomorphism $\psi:\Pi\to S$ with $\psi|_J=\sigma$, $\psi(T)=w_{12}$, and $\phi_n\psi=p$. Its image contains every generator $X_{ij}(c)$ with $i,j\geq2$, and every generator $X_{1j}(c)$. The conjugate of $X_{i2}(c)$ by $w_{12}$ is $X_{i1}(c)$ for $i\geq3$, while the conjugate of $X_{12}(c)$ is $X_{21}(c)$. These are all the defining generators of $S$, so $\psi$ is surjective. Since $p$ is an isomorphism, $\psi p^{-1}:G\to S$ is a surjective section of $\phi_n$, and is therefore its inverse.
\end{proof}

To apply the criterion over the binary Leavitt ring, we need two algebraic facts. The next lemma uses an additional index, as in \cite[Chapter~III, proof of Theorem~5.2.1]{Weibel13}. We include the proof because centrality is asserted after one stabilization.

\begin{lemma}
\label{lem:padded-centrality}
For any unital ring $B$ and $m\geq3$, $j_m(N_m(B))$ is central in $\St_{m+1}(B)$.
\end{lemma}

\begin{proof}
For a column $a\in B^m$ and row $b\in M_{1,m}(B)$ put $u(a)=\prod_{i=1}^mX_{i,m+1}(a_i)$ and $v(b)=\prod_{i=1}^mX_{m+1,i}(b_i)$. The factors in each product commute. The restriction of $\phi_{m+1}$ to each of these subgroups is injective, since its matrix image determines every coefficient. For $i,j\leq m$, the Steinberg relations imply
$$
X_{ij}(c)u(a)X_{ij}(c)^{-1}=u(t_{ij}(c)a)\qquad\text{and}\qquad X_{ij}(c)v(b)X_{ij}(c)^{-1}=v(bt_{ij}(c)^{-1}).
$$
Only $a_i$ changes in the first formula, by addition of $ca_j$; only $b_j$ changes in the second, by addition of $-b_i c$. Thus a word on the first $m$ indices acts on these subgroups through its matrix image, and every element of $j_m(N_m(B))$ centralizes both subgroups. Each remaining generator is $X_{ij}(c)=[X_{i,m+1}(c),X_{m+1,j}(1)]$, so it too is centralized. This proves the assertion.
\end{proof}

We deduce the following equality from the $GE$ result of Menal--Moncasi and Ara--Goodearl--Pardo, together with the acyclicity already established.

\begin{lemma}
\label{lem:elementary-generation}
For $R=L_{\F_2}(1,2)$, one has $\GL_m(R)=E_m(R)$ for every $m\geq2$.
\end{lemma}

\begin{proof}
The algebra $R$ is $V_{1,2}$ in \cite[\S4]{AGP02}: the universal row is $(e,f)$ and the inverse column is $(e^*,f^*)^t$. Theorem~4.2 of that paper shows that $R$ is purely infinite simple. The proof of its Theorem~2.3 invokes Menal--Moncasi's $GE$ argument \cite[Theorem~2.2 and the remark after Corollary~2.3]{MM84}. The latter remark applies to a simple ring in which each nonzero $x$ admits $y,z$ with $yxz=1$, and requires no regularity assumption. Thus every invertible matrix over $R$ is a product of elementary matrices and invertible diagonal matrices.

By Theorem~\ref{thm:acyclicity}, the group $R^\times$ is perfect. We show directly that each diagonal factor belongs to $E_m(R)$ for every $m\geq2$. In a block on two coordinates, put $W(u)=t_{12}(u)t_{21}(-u^{-1})t_{12}(u)$. A direct multiplication shows that $W(u)W(-1)=\operatorname{diag}(u,u^{-1})$. For units $a,b$,
$$
\operatorname{diag}([a,b],1) =\operatorname{diag}(a,a^{-1}) \operatorname{diag}(b,b^{-1}) \operatorname{diag}(a^{-1}b^{-1},ba).
$$
Every factor on the right belongs to $E_2(R)$. Every unit is a finite product of commutators, so $\operatorname{diag}(u,1)\in E_2(R)$ for every $u\in R^\times$. Permuting indices and using one other coordinate proves the same assertion in any chosen diagonal position. Thus every invertible diagonal matrix belongs to $E_m(R)$, and the $GE$ result proves the lemma.
\end{proof}

\begin{theorem}
\label{thm:steinberg-isomorphism}
For $R=L_{\F_2}(1,2)$ and every $r\geq3$, the canonical map $\phi_r:\St_r(R)\to\GL_r(R)$ is an isomorphism.
\end{theorem}

\begin{proof}
Write $S_r=\St_r(R)$, $G_r=\GL_r(R)$, and $N_r=\ker\phi_r$. By Lemma~\ref{lem:elementary-generation}, $\phi_3$ is surjective. The group $S_3$ is perfect, since each generator is a commutator using a third index. The five-term integral homology sequence for $1\to N_3\to S_3\to G_3\to1$ contains
$$
0=H_2(G_3;\Z)\longrightarrow N_3/[S_3,N_3] \longrightarrow H_1(S_3;\Z)=0,
$$
where the first equality follows from Theorem~\ref{thm:acyclicity} and the coordinate isomorphisms of Lemma~\ref{lem:leaf-coordinates}. Consequently $N_3=[S_3,N_3]$. In view of Lemma~\ref{lem:padded-centrality}, $j_3(N_3)=[j_3(S_3),j_3(N_3)]=1$ in $S_4$.

The complex $X_4(R)$ is the ordered frame complex used in Proposition~\ref{prop:frame-cone}, with complements forgotten. According to that proposition, any finite collection of vertices and edges has a common apex, with complements for all the resulting frames. Coning a pair of vertices produces a path between them. For a finite edge loop, choose one such apex $v$. For each ordered edge $(u,w)$ traversed by the loop, use the triangle $(v,u,w)$. Adjacent triangles share the radial edge $(v,u)$, so they glue to a continuous triangular fan whose boundary is the prescribed loop. A reverse traversal uses the same ordered edge and its triangle with reverse orientation; it does not replace that edge by the distinct cell $(w,u)$. Repeated vertices cause no difficulty for this map of a disk. By cellular approximation, every element of the fundamental group is represented by a finite edge loop, and the fan contracts it. Thus $|X_4(R)|$ is connected and simply connected.

By Lemma~\ref{lem:elementary-generation}, $G_2=E_2(R)$ and $G_3=E_3(R)$. All hypotheses of Theorem~\ref{thm:raw-criterion} are therefore satisfied at $n=4$, so $\phi_4$ is an isomorphism.

Coordinate refinement transfers this conclusion to the other ranks. Let $T_r:R^r\to R^{r+1}$ fix the first $r-1$ coordinates and send the last coordinate $x$ to $(e^*x,f^*x)^t$. Let $U_r:R^{r+1}\to R^r$ fix the first $r-1$ coordinates and send the last pair $(y,z)$ to $ey+fz$. By the Leavitt identities, $U_rT_r=I_r$ and $T_rU_r=I_{r+1}$. Hence $\alpha_r(A)=T_rAU_r$ is a unital isomorphism of matrix rings.

For $r\geq3$, this is covered by the Steinberg isomorphism $D_r:S_r\to S_{r+1}$ with formulas
$$
\begin{aligned}
\eqrownum &\quad D_rX_{ij}(a)=X_{ij}(a) \quad \text{for }i,j<r,\\
\eqrownum &\quad D_rX_{ir}(a)=X_{ir}(ae)X_{i,r+1}(af) \quad \text{for }i<r,\\
\eqrownum[eq:steinberg-refinement] &\quad D_rX_{rj}(a)=X_{rj}(e^*a)X_{r+1,j}(f^*a) \quad \text{for }j<r.
\end{aligned}
$$
To justify this assertion, apply \cite[\S4, Proposition~1]{Voronetsky20} to the ordinary groups, taking $S=\{1\}$ as explained immediately before Lemma~5, for the ring $M_{r+1}(R)$. Its complete family of standard orthogonal idempotents is full and Morita equivalent, as witnessed by $E_{aa}=E_{ab}E_{bb}E_{ba}$. Merge the last two idempotents. The maps $T_r,U_r$ identify the resulting coarse presentation with $S_r$. The proposition identifies it with the fine presentation $S_{r+1}$ whenever the fine family has at least four members; here that number is $r+1\geq4$. Its map on generators is precisely \eqref{eq:steinberg-refinement}. This application requires no hypothesis on stable rank, quasi-finiteness, or centrality of the kernel.

The maps satisfy $\phi_{r+1}D_r=\alpha_r\phi_r$, and both $D_r$ and $\alpha_r$ are isomorphisms. Therefore $D_r(N_r)=N_{r+1}$. Starting from $N_4=1$ and iterating in either direction proves $N_r=1$ for every $r\geq3$. Surjectivity follows from Lemma~\ref{lem:elementary-generation}, and completes the proof.
\end{proof}

\section{An explicit finite presentation}
\label{sec:finite-presentation}

Finite presentability is a consequence of Theorem~\ref{thm:steinberg-isomorphism} and Krsti\'c--McCool's theorem \cite[Theorem~3]{KM99}. The relation of this question to unstable $K_2$ was already considered in \cite[\S8]{KT26}. We first record the ring quotient and then specialize the bounded construction of \cite[\S3]{KM99}, including its finite index sets and the evaluation of the generators in $R^\times$.

\begin{theorem}\label{thm:finite-presentation}
The group $G=R^\times$ is finitely presented.
\end{theorem}

\begin{proof}
Let $F=\Z\langle e,f,s,t\rangle$ and let $J$ be the two-sided ideal generated by
$$
2,\quad se-1,\quad tf-1,\quad sf,\quad te,\quad es+ft-1. \eqnum[eq:ring-ideal-generators]
$$
There is a unital ring isomorphism $F/J\cong R$ sending $s$ to $e^*$ and $t$ to $f^*$. Indeed, the displayed relations hold in $R$, while $2=0$ makes $F/J$ an $\F_2$-algebra whose generators satisfy the Leavitt relations. The two maps supplied by these universal properties fix the four generators and are inverse.

Krsti\'c--McCool \cite[Theorem~3]{KM99} prove that $\St_n(B)$ is finitely presented for every finitely presented associative unital ring $B$ and every $n\geq4$. Their definition of a finitely presented ring is a quotient of a free associative unital $\Z$-algebra on finitely many generators by a finitely generated two-sided ideal. Thus $\St_5(R)$ is finitely presented. By Theorem~\ref{thm:steinberg-isomorphism}, $\phi_5\colon\St_5(R)\to\GL_5(R)$ is an isomorphism. The coordinates from the five leaves below identify $\GL_5(R)$ with $G$.
\end{proof}

The following argument for the quotient ring, used in \cite[pp.~178--179]{KM99}, determines the additional relators. For $n\geq3$, let $K$ be the normal closure in $\St_n(F)$ of the elements $X_{ij}(y)$, where $i\ne j$ and $y$ ranges over \eqref{eq:ring-ideal-generators}. Put
$$
I=\{a\in F:X_{ij}(a)\in K\text{ for every }i\ne j\}.
$$
The additivity relation makes $I$ an additive subgroup. For $a\in I$, $b\in F$ and $k\notin\{i,j\}$, the identities
$$
X_{ij}(ba)=[X_{ik}(b),X_{kj}(a)]\qquad\text{and}\qquad X_{ij}(ab)=[X_{ik}(a),X_{kj}(b)]
$$
and normality of $K$ show that $ba,ab\in I$. Hence $J\subseteq I$. Conversely, every normal generator of $K$ maps to the identity in $\St_n(F/J)$. In $\St_n(F)/K$, the value of $X_{ij}(a)$ therefore depends only on $a+J$: if $a-b\in J$, then $X_{ij}(a)X_{ij}(b)^{-1}=X_{ij}(a-b)\in K$ by the additivity relation. These elements satisfy all Steinberg relations and define an inverse to the natural quotient map. Consequently
$$
\St_n(F)/K\cong\St_n(F/J). \eqnum[eq:steinberg-ring-quotient]
$$
This proves both inclusions in the assertion about the kernel.

We specialize the presentation of Krsti\'c and McCool \cite[\S3]{KM99} to the four ring generators and six ring relations above. Let $\mathcal T=\{e,f,s,t\}$ be an alphabet, let $\mathcal T^*$ be its free monoid, and denote the empty word by $\varepsilon$. All words in the presentation are literal words in this alphabet; their multiplication is concatenation. Set $m=194$ and
$$
\mathcal W=\{u\in\mathcal T^*:1\leq |u|\leq m\}\qquad\text{and}\qquad \mathcal W^0=\mathcal W\cup\{\varepsilon\}.
$$
For $d=2,3,4$, let $\mathcal I_d$ be the set of ordered $d$-tuples of pairwise distinct elements of $\{1,2,3,4,5\}$. Define the finite coefficient sets
$$
\begin{aligned}
\mathcal A&=\{(u,v)\in\mathcal W^2:|u|+|v|\leq m\},\\
 \mathcal B&=\{(u,v)\in\mathcal W^2:|u|+|v|\leq m+3\},\\
 \mathcal C&=\{(u,v)\in\mathcal W^2:|u|<m\}\qquad\text{and}\qquad\mathcal D=\mathcal C,\\
 \mathcal P&=\{(\varepsilon,\varepsilon)\}
       \cup\{(\varepsilon,a),(a,\varepsilon):a\in\mathcal T\}.
\end{aligned}
$$
In particular, both words in $\mathcal B$ have length at most $m$, and the second word in $\mathcal C$ or $\mathcal D$ has length at most $m$. The strict inequality in the first coefficient is required in both sets. The set $\mathcal P$ has nine ordered pairs.

Take the generators
$$
x_{ij}(u)\qquad\text{for }(i,j)\in\mathcal I_2\text{ and }u\in\mathcal W^0. \eqnum[eq:finite-generators]
$$
The symbol $x_{ij}(\varepsilon)$ is a generator, not the group identity. For the following five families, each relation is imposed for every index tuple and coefficient pair in the indicated finite sets:
$$
\begin{aligned}
\eqrowtag{A}{eq:finite-A} &\quad [x_{ij}(u),x_{jk}(v)]=x_{ik}(uv)\\
&\quad \text{for }(i,j,k)\in\mathcal I_3\text{ and }(u,v)\in\mathcal A\cup\mathcal P;\\
\eqrowtag{B}{eq:finite-B} &\quad [x_{ij}(u),x_{kl}(v)]=1\\
&\quad \text{for }(i,j,k,l)\in\mathcal I_4\text{ and }(u,v)\in\mathcal B\cup\mathcal P;\\
\eqrowtag{$C_\ell$}{eq:finite-C-left} &\quad [x_{ik}(u),x_{jk}(v)]=1\\
&\quad \text{for }(i,j,k)\in\mathcal I_3\text{ and }(u,v)\in\mathcal C\cup\mathcal P;\\
\eqrowtag{$C_r$}{eq:finite-C-right} &\quad [x_{ki}(u),x_{kj}(v)]=1\\
&\quad \text{for }(i,j,k)\in\mathcal I_3\text{ and }(u,v)\in\mathcal C\cup\mathcal P;\\
\eqrowtag{D}{eq:finite-D} &\quad [x_{ij}(u),x_{ij}(v)]=1\\
&\quad \text{for }(i,j)\in\mathcal I_2\text{ and }(u,v)\in\mathcal D\cup\mathcal P.
\end{aligned}
$$
Thus every family includes $(\varepsilon,\varepsilon)$ and both $(\varepsilon,a)$ and $(a,\varepsilon)$ for each letter $a$. In \eqref{eq:finite-A}, the word $uv$ has length at most $m$, or at most one for a pair in $\mathcal P$, so its generator is defined.

For each $(i,j)\in\mathcal I_2$, impose also the following six relations:
$$
\begin{aligned}
\eqrowtag{$L_1$}{eq:finite-L1} &\quad x_{ij}(\varepsilon)^2=1,\\
\eqrowtag{$L_2$}{eq:finite-L2} &\quad x_{ij}(se)x_{ij}(\varepsilon)^{-1}=1,\\
\eqrowtag{$L_3$}{eq:finite-L3} &\quad x_{ij}(tf)x_{ij}(\varepsilon)^{-1}=1,\\
\eqrowtag{$L_4$}{eq:finite-L4} &\quad x_{ij}(sf)=1,\\
\eqrowtag{$L_5$}{eq:finite-L5} &\quad x_{ij}(te)=1,\\
\eqrowtag{$L_6$}{eq:finite-L6} &\quad x_{ij}(es)x_{ij}(ft)x_{ij}(\varepsilon)^{-1}=1.
\end{aligned}
$$
These are $120$ labelled relations. The inverse signs are retained because the presentation preceding them is over $\Z$; characteristic two is imposed by \eqref{eq:finite-L1}.

\begin{proposition}\label{prop:explicit-presentation}
Let $\Pi$ be the group with generators \eqref{eq:finite-generators} and relations \eqref{eq:finite-A}--\eqref{eq:finite-D} and \eqref{eq:finite-L1}--\eqref{eq:finite-L6}, with exactly the finite index sets specified above. Put
$$
(\alpha_1,\alpha_2,\alpha_3,\alpha_4,\alpha_5) =(e,fe,f^2e,f^3e,f^4).
$$
For $u\in\mathcal T^*$, write $\bar u$ for its value in $R$ under $s\mapsto e^*$, $t\mapsto f^*$ and $\varepsilon\mapsto1$. Then
$$
x_{ij}(u)\longmapsto 1+\alpha_i\bar u\alpha_j^* \eqnum[eq:finite-evaluation]
$$
extends to an isomorphism $\Pi\cong G$.
\end{proposition}

\begin{proof}
The bounded construction in \cite[Section~3, pp.~179--182]{KM99} identifies the group given by \eqref{eq:finite-A}--\eqref{eq:finite-D} with $\St_5(F)$, taking $x_{ij}(u)$ to $X_{ij}(u)$. Its sufficient bound is $m>3|\mathcal T|^3+1$, which holds since $194>193$. The relations indexed by nonempty words in $\mathcal T^*$ use the sets $\mathcal A$, $\mathcal B$, $\mathcal C$ and $\mathcal D$ above. The generators $x_{ij}(\varepsilon)$ are adjoined with all instances of these four types having an empty coefficient and total length at most one, as indexed by $\mathcal P$; the induction in that proof supplies the remaining constant relations.

In $\St_5(F)$, the additivity relation identifies \eqref{eq:finite-L1}--\eqref{eq:finite-L6} with $X_{ij}(y)=1$ for the six polynomials in \eqref{eq:ring-ideal-generators}. It follows from \eqref{eq:steinberg-ring-quotient} that there is an isomorphism
$$
\kappa\colon\Pi\xrightarrow{\ \cong\ }\St_5(R),\qquad x_{ij}(u)\longmapsto X_{ij}(\bar u).
$$

Let $L=(\alpha_1,\ldots,\alpha_5)$ and let $L^*$ be the column of the words $\alpha_i^*$. Since no distinct leaves are comparable by prefix, $\alpha_i^*\alpha_j=\delta_{ij}$. Completeness follows by repeatedly substituting $ff^*=1-ee^*$ in
$$
\sum_{a=0}^3 f^aee^*(f^*)^a+f^4(f^*)^4=1.
$$
Thus $LL^*=1$ and $L^*L=I_5$. The mutually inverse unital ring maps are
$$
\Theta(A)=LAL^*\qquad\text{and}\qquad \Theta^{-1}(a)=L^*aL=(\alpha_i^*a\alpha_j)_{i,j}. \eqnum[eq:five-leaf-coordinates]
$$
For example, multiplication of images inserts $L^*L=I_5$, and the two inverse identities follow from these same two leaf identities.

The composite $\Theta\phi_5\kappa$ sends each generator to \eqref{eq:finite-evaluation}. All three maps are isomorphisms, the middle one by Theorem~\ref{thm:steinberg-isomorphism}. This proves the assertion. The displayed generator image has two-sided inverse $1-\alpha_i\bar u\alpha_j^*$, since $(\alpha_i\bar u\alpha_j^*)^2=0$ for $i\ne j$. In $R$ this inverse equals the generator image itself.
\end{proof}

One can also describe the inverse of \eqref{eq:finite-evaluation} in words. Given $g\in G$, form the matrix $A=(\alpha_i^*g\alpha_j)_{i,j}$ and choose an elementary factorization $A=\prod_{\nu=1}^q(I_5+a_\nu E_{i_\nu j_\nu})$. Such a factorization exists since $\GL_5(R)=E_5(R)$. Lift each $a_\nu$ to a finite integer linear combination of words in $F$. Represent a coefficient sum by a product in the corresponding root group, and an integer multiple by a power. For a word $u$ of length at most $m$, use the generator $x_{ij}(u)$; the empty word uses $x_{ij}(\varepsilon)$. For a longer word $u=va$, with $a$ its last letter, recursively use $[x_{ik}(v),x_{kj}(a)]$, where $k$ is the least index outside $\{i,j\}$ and the first factor is expanded by the same rule. The resulting product represents the inverse image of $g$ in $\Pi$. Its independence of the polynomial lifts and the elementary factorization follows from the isomorphisms just proved.

The presentation has $20(1+4+\cdots+4^{194})=20(4^{195}-1)/3$ generators. Every generator and relator is indexed by one of the finite sets specified above.

\clearpage
\begin{thebibliography}{99}

\bibitem{AAP08}
G.~Abrams, P.~N.~\'{A}nh and E.~Pardo, \emph{Isomorphisms between Leavitt algebras and their matrix rings}, J. Reine Angew. Math. \textbf{624} (2008), 103--132. \url{https://doi.org/10.1515/CRELLE.2008.082}.

\bibitem{AAS17}
G.~Abrams, P.~Ara and M.~Siles Molina, \emph{Leavitt Path Algebras}, Lecture Notes in Mathematics, vol.~2191, Springer, London, 2017. \url{https://doi.org/10.1007/978-1-4471-7344-1}.

\bibitem{ABC09}
P.~Ara, M.~Brustenga and G.~Corti\~nas, \emph{$K$-theory of Leavitt path algebras}, M\"unster J. Math. \textbf{2} (2009), 5--34. \url{https://arxiv.org/abs/0903.0056v2}.

\bibitem{AGP02}
P.~Ara, K.~R.~Goodearl and E.~Pardo, \emph{$K_0$ of purely infinite simple regular rings}, $K$-Theory \textbf{26} (2002), no.~1, 69--100. \url{https://doi.org/10.1023/A:1016358107918}.

\bibitem{Brown84}
K.~S.~Brown, \emph{Presentations for groups acting on simply-connected complexes}, J. Pure Appl. Algebra \textbf{32} (1984), 1--10. \url{https://doi.org/10.1016/0022-4049(84)90009-4}.

\bibitem{Freeland19}
R.~L.~Freeland, \emph{Relating Thompson's group $V$ to graphs of groups and Hecke algebras}, Ph.D. thesis, University of Cambridge, 2019. \url{https://doi.org/10.17863/CAM.52134}.

\bibitem{Gorazd26}
R.~Gorazd, \emph{Embedding Higman--Thompson groups of unfolding trees into the Leavitt path algebras}, Indag. Math. (2026), in press. \url{https://doi.org/10.1016/j.indag.2026.03.001}.

\bibitem{KT26}
H.~V.~Khanh and V.~H.~Thanh, \emph{Matrix generators for the unit groups of $L_K(1,d)$}, arXiv:2607.10351v1 (2026). \url{https://arxiv.org/abs/2607.10351v1}.

\bibitem{Knudson02}
K.~P.~Knudson, \emph{Unstable homotopy invariance for finite fields}, Fund. Math. \textbf{175} (2002), no.~2, 155--162. \url{https://doi.org/10.4064/fm175-2-5}.

\bibitem{KM99}
S.~Krsti\'c and J.~McCool, \emph{Presenting $\mathrm{GL}_n(k\langle T\rangle)$}, J. Pure Appl. Algebra \textbf{141} (1999), no.~2, 175--183. \url{https://doi.org/10.1016/S0022-4049(98)00022-X}.

\bibitem{Li25}
X.~Li, \emph{Ample groupoids, topological full groups, algebraic $K$-theory spectra and infinite loop spaces}, Forum Math. Pi \textbf{13} (2025), article~e9. \url{https://doi.org/10.1017/fmp.2024.31}.

\bibitem{MM84}
P.~Menal and J.~Moncasi, \emph{$K_1$ of von Neumann regular rings}, J. Pure Appl. Algebra \textbf{33} (1984), no.~3, 295--312. \url{https://doi.org/10.1016/0022-4049(84)90064-1}.

\bibitem{NS90}
Yu.~P.~Nesterenko and A.~A.~Suslin, \emph{Homology of the full linear group over a local ring, and Milnor's $K$-theory}, Math. USSR-Izv. \textbf{34} (1990), no.~1, 121--145. \url{https://doi.org/10.1070/IM1990v034n01ABEH000610}.

\bibitem{Pardo11}
E.~Pardo, \emph{The isomorphism problem for Higman--Thompson groups}, J. Algebra \textbf{344} (2011), 172--183. \url{https://doi.org/10.1016/j.jalgebra.2011.07.026}.

\bibitem{PW25}
M.~Palmer and X.~Wu, \emph{Embedding groups into acyclic groups}, arXiv:2510.16879v1 (2025). \url{https://arxiv.org/abs/2510.16879v1}.

\bibitem{Preusser26}
R.~Preusser, \emph{The general linear groupoid of a Leavitt path algebra}, arXiv:2608.21026v1 (2026). \url{https://arxiv.org/abs/2608.21026v1}.

\bibitem{Quillen72}
D.~Quillen, \emph{On the cohomology and $K$-theory of the general linear groups over a finite field}, Ann. of Math. (2) \textbf{96} (1972), no.~3, 552--586. \url{https://doi.org/10.2307/1970825}.

\bibitem{RWW17}
O.~Randal-Williams and N.~Wahl, \emph{Homological stability for automorphism groups}, Adv. Math. \textbf{318} (2017), 534--626. \url{https://doi.org/10.1016/j.aim.2017.07.022}.

\bibitem{Schlichting17}
M.~Schlichting, \emph{Euler class groups and the homology of elementary and special linear groups}, Adv. Math. \textbf{320} (2017), 1--81. \url{https://doi.org/10.1016/j.aim.2017.08.034}.

\bibitem{Suslin84}
A.~A.~Suslin, \emph{Homology of $\mathrm{GL}_n$, characteristic classes and Milnor $K$-theory}, Trudy Mat. Inst. Steklov. \textbf{165} (1984), 188--204 (Russian); English translation, Proc. Steklov Inst. Math. \textbf{165} (1985), 207--226. \url{https://www.mathnet.ru/eng/tm2280}.

\bibitem{SW19}
M.~Szymik and N.~Wahl, \emph{The homology of the Higman--Thompson groups}, Invent. Math. \textbf{216} (2019), 445--518. \url{https://doi.org/10.1007/s00222-018-00848-z}. Preprint version: \url{https://arxiv.org/abs/1411.5035v5}.

\bibitem{vdK80}
W.~van der Kallen, \emph{Homology stability for linear groups}, Invent. Math. \textbf{60} (1980), no.~3, 269--295. \url{https://doi.org/10.1007/BF01390018}.

\bibitem{Voronetsky20}
E.~Voronetsky, \emph{Centrality of $K_2$-functor revisited}, J. Pure Appl. Algebra \textbf{225} (2021), no.~4, article~106547. \url{https://doi.org/10.1016/j.jpaa.2020.106547}. The version used here is \url{https://arxiv.org/abs/2004.08551v2}.

\bibitem{Wagoner72}
J.~B.~Wagoner, \emph{Delooping classifying spaces in algebraic $K$-theory}, Topology \textbf{11} (1972), 349--370. \url{https://doi.org/10.1016/0040-9383(72)90031-6}.

\bibitem{Weibel94}
C.~A.~Weibel, \emph{An Introduction to Homological Algebra}, Cambridge Studies in Advanced Mathematics, vol.~38, Cambridge University Press, Cambridge, 1994. \url{https://doi.org/10.1017/CBO9781139644136}.

\bibitem{Weibel13}
C.~A.~Weibel, \emph{The $K$-book: An Introduction to Algebraic $K$-theory}, Graduate Studies in Mathematics, vol.~145, American Mathematical Society, Providence, RI, 2013. \url{https://doi.org/10.1090/gsm/145}.

\end{thebibliography}

\end{document}
